\documentclass[a4paper]{article}

%%%%%%%% CREATE DOCUMENT STRUCTURE %%%%%%%%
%% Language and font encodings
\usepackage[spanish]{babel}
\usepackage[utf8x]{inputenc}
\usepackage[T1]{fontenc}

\usepackage{listings} % Para incluir código

\usepackage{xcolor}   % Colores para el código
% Definición de colores
\definecolor{codegreen}{rgb}{0,0.6,0}
\definecolor{codegray}{rgb}{0.5,0.5,0.5}
\definecolor{codepurple}{rgb}{0.58,0,0.82}
\definecolor{backcolour}{rgb}{0.95,0.95,0.92}
\definecolor{bluekeywords}{rgb}{0.13,0.13,1}

% Configuración del estilo para PostScript
\lstdefinelanguage{PostScript}{
  morekeywords={
    moveto, lineto, rlineto, arc, closepath, stroke, fill, 
    newpath, showpage, setlinewidth, setgray, setrgbcolor,
    findfont, scalefont, setfont, show, def, setpagedevice
  },
  sensitive=true,
  morecomment=[l]{\%},
  %morestring=[b](,
  %morestring=[b][\color{codepurple}]{)}
  stringstyle=\color{codepurple}
}

\lstset{        
    backgroundcolor=\color{backcolour},
    basicstyle=\ttfamily\small, % Fuente monoespaciada
    %breakatwhitespace=false,         
    breaklines=true,    % Saltar líneas largas      
    captionpos=b,   
    commentstyle=\color{codegreen},
    extendedchars=true, % Para las tildes
    frame=single,               % Marco alrededor del código
    keepspaces=true,                 
    keywordstyle=\color{bluekeywords}\bfseries,
    language=PostScript,
    literate= {á}{{\'a}}1 {é}{{\'e}}1 {í}{{\'i}}1 {ó}{{\'o}}1 {ú}{{\'u}}1 {ñ}{{\~n}}1 {Á}{{\'A}}1 {É}{{\'E}}1 {Í}{{\'I}}1 {Ó}{{\'O}}1 {Ú}{{\'U}}1 {Ñ}{{\~N}}1, % Para las tildes y caracteres del español
    numbers=left,   % Números de línea a la izquierda                 
    numbersep=5pt,   
    numberstyle=\tiny\color{codegray},           
    showspaces=false,                
    showstringspaces=false,
    showtabs=false,     
    stringstyle=\color{codepurple},                  
    tabsize=2
}

%% Sets page size and margins
\usepackage[a4paper,top=3cm,bottom=2cm,left=2cm,right=2cm,marginparwidth=1.75cm]{geometry}

%% Useful packages
\usepackage{amsmath}
\usepackage{graphicx}
\usepackage[colorinlistoftodos]{todonotes}
%\usepackage[colorlinks=true, allcolors=blue]{hyperref}
\usepackage{caption}
\usepackage{subcaption}
\usepackage{sectsty}

\usepackage{float}
\usepackage{titling} 
\usepackage{movie15} %Para añadir video

%% Referencias/Bibliografía
%\usepackage[backend=biber, style=apa]{biblatex}
%\addbibresource{ref.bib} % Nombre del archivo con las referencias


\usepackage[colorinlistoftodos]{todonotes}
\definecolor{darkgreen}{rgb}{0.0, 0.4, 0.0}


%Links de referencias
\usepackage[hidelinks]{hyperref} % Hace que las referencias sean clicables
\newcommand{\myref}[1]{\textsuperscript{\textbf{\ref{#1}}}} % Define un estilo propio

\usepackage{pdfpages} % Para añadir PDFs (vista esquemática)


\usepackage{pdflscape} % Para páginas en orientación horizontal

%%%%%%%% DOCUMENT %%%%%%%%
\begin{document}

%%%% Title Page
\begin{titlepage}

\newcommand{\HRule}{\rule{\linewidth}{0.5mm}} 							% horizontal line and its thickness
\center 
 
% University
\textsc{\LARGE Universidad de Granada}\\[1cm]

% Document info
\textsc{\Large Periféricos y Dispositivos de Interfaz Humana}\\[0.2cm]
\textsc{\large Curso 2025/2026}\\[1cm] 										% Course Code
\HRule \\[0.8cm]
{ \huge \bfseries Práctica 4. El lenguaje Postscript}\\[0.7cm]								% Assignment
\HRule \\[2cm]
\large
\emph{Autora:}\\
Inés Ruiz Sánchez\\[1.5cm]													% Author info
{\large \today}\\[5cm]

\includegraphics[width=0.3\textwidth]{images/UGRlogo.png}\\[1cm] 	% University logo
\vfill 
\end{titlepage}

%%\begin{abstract}
%%Your abstract.
%%\end{abstract}

% Generar índice de contenidos y figuras
\tableofcontents
%\newpage
\listoffigures
\newpage

%%%% SECTIONS
\section{Introducción}
El objetivo principal es conocer cómo utilizar el lenguaje Postscript a nivel básico. \\

Los objetivos concretos de esta práctica son: 
\begin{itemize}
    \item Conocer la sintaxis básica del lenguaje Postscript.
    \item Crear varios archivos y probar el funcionamiento con un visor Postscript. 
    \item Transformar los ejemplos creados a otros formatos de archivo, como PDF. 
\end{itemize}

%\textit{Nota}: Las tildes en los trozos de código no se reconocen en \textit{LATEX} por lo que el carácter no se escribe. Es por eso por lo que o bien no aparece un caracter o bien el caracter no tiene tilde, pareciendo una falta ortográfica.\\

En el anexo\myref{anexo:an} se incluyen los ficheros de códigos al completo. Para obtener los documentos PDFs de esta práctica he usado la página \url{https://online2pdf.com/en/ps2pdf} recomendada por el profesor.\\

Repositorio en GitHub: \url{https://github.com/ruiz314/PDIH/tree/main/P4}

\section{Requisitos mínimos}
Para superar esta práctica los dos proyectos que deberán realizarse serán los siguientes: 

\begin{itemize}
    \item Página en la que se muestre el dibujo de unos árboles, similar al siguiente\myref{fig:1ejemplo1}:
    \begin{figure}[H]
        \includegraphics[width=0.7\textwidth]{images/ejemplos_guion/ejemplo1.png}
        \caption{Ejercicio 1. Árboles}
        \label{fig:1ejemplo1}
    \end{figure}
    
    \item Página en la que se muestre el diseño de una tarjeta de visita, similar a la siguiente\myref{fig:2ejemplo2}:
    \begin{figure}[H]
        \includegraphics[width=0.7\textwidth]{images/ejemplos_guion/ejemplo2.png}
        \caption{Ejercicio 2. Tarjeta de visita}
        \label{fig:2ejemplo2}
    \end{figure}
    
    \item Página en la que se muestre el dibujo de un corazón y una poesía\myref{fig:3ejemplo3}:
    \begin{figure}[H]
        \includegraphics[width=0.7\textwidth]{images/ejemplos_guion/ejemplo3.png}
        \caption{Ejercicio 3. Dibujo con poesía}
        \label{fig:3ejemplo3}
    \end{figure}
\end{itemize}

\subsection{Ejercicio 1}\label{subsec:ej1}
%\subsection{Ejercicio 1}\label{subsec:ej1}

En este primer ejercicio se ha creado una escena que incluye dos tipos de árboles (un árbol de copa redonda y un pino), una pelota, el suelo y un sol.

Para no tener problemas con la superposición en los árboles primero se dibuja el tronco, y a continuación las hojas. De esta forma las hojas tapan un poco al tronco, dando realismo. Para dibujar el árbol de la izquierda se han usado varios círculos rellenos, con distinta posición y distinto radio. Para ello se ha usado la función \textit{arc}. \\

Para el pino (árbol a la derecha) se han usado triángulos de distinto tamaño:
\begin{lstlisting}[language=PostScript, caption=Pino, captionpos=b, label=fun:pino]
% Hojas de abajo
newpath
    0.3 0.7 0 setrgbcolor  % Color verde
    285 90 moveto         
    230 0 rlineto
    -115 310 rlineto
    closepath
    fill
stroke

% Hojas de medio-bajo
newpath
    0.3 0.7 0 setrgbcolor  % Color verde
    300 150 moveto         
    200 0 rlineto
    -100 250 rlineto
    closepath
    fill
stroke

% Hojas del medio-alto
newpath
    0.3 0.7 0 setrgbcolor  % Color verde
    315 210 moveto         
    170 0 rlineto
    -85 190 rlineto
    closepath
    fill
stroke

% Hojas de arriba
newpath
    0.3 0.7 0 setrgbcolor  % Color verde
    330 270 moveto         
    140 0 rlineto
    -70 130 rlineto
    closepath
    fill
stroke
\end{lstlisting}
Se puede observar que para cada triángulo la coordenada $Y$ es mayor, para que esté situado un poco más arriba que el anterior.\\

También se han usado distintos colores para las hojas del árbol de copa redonda para dar sensación de luz, sombras y hojas nuevas. Se ha aplicado color a todos los elementos con \textit{setrgbcolor}, mezclando valores del `0' al `1' para los canales Rojo, Verde y Azul.\\

En la esquina superior derecha se ha dibujado un \textbf{sol con rayos}:
\begin{lstlisting}[language=PostScript, caption=Sol, captionpos=b, label=fun:sol]
newpath
    1 0.83 0 setrgbcolor  % Color amarillo-naranja
    600 750 40 0 360 arc % Pintar círculo cerca de la esquina superior derecha
stroke
\end{lstlisting}
Para dibujar el círculo del sol se usa la instrucción \textit{arc} pero no se ve entero porque se ha colocado en el límite ($x=600$). Otra opción sería cambiar el prámetro $360$ por los grados concretos de la circunferencia.\\

Para pintar los rayos de luz del sol se pintan tres rayas con un grosor \textit{5 setlinewidth} porque es el que se ha dado a la pelota y no se ha modificado posteriormente. Se sitúa el ``lápiz'' en $X = 550$ y se mueve hacia la izquierda $80$ unidades. Como se han dibujado 3 rayos, se necesitan tres bloques que cambien el valor de la componente $Y$:
\begin{lstlisting}[language=PostScript, caption=Rayos del sol, captionpos=b, label=fun:rayos_sol]
% Rayos del SOL
newpath % Rayo horizontal
    1 0.83 0 setrgbcolor  % Color amarillo-naranja
    550 750 moveto
    -80 0 rlineto
stroke

newpath % Rayo oblicuo arriba
    1 0.83 0 setrgbcolor  % Color amarillo-naranja
    550 780 moveto
    -80 25 rlineto
stroke

newpath % Rayo oblicuo abajo
    1 0.83 0 setrgbcolor  % Color amarillo-naranja
    550 720 moveto
    -80 -25 rlineto
stroke
\end{lstlisting}

A continuación se muestra la imagen obtenida\myref{fig:1ejer1}:
\begin{figure}[H]
\includegraphics[width=0.8\textwidth]{images/ejer1_completo.jpg}
\caption{Propuesta ejercicio 1}
\label{fig:1ejer1}
\end{figure}

Como \textbf{ampliación} para los árboles se podrían dibujar las ramas que sujetan las hojas en el de copa redonda. O dibujar algún fruto como unas manzanas (círculos rojos). Para el pino se podría haber decorado como un pino de Navidad. Otra opción de mejora sería darle realismo a la pelota haciéndole los dibujos de una pelota de tenis o baloncesto. Por último, se podrían incluir algunas nubes o pájaros en el cielo.\\
 
El código completo se puede ver en el anexo \ref{anexo:ej1}. En el anexo \ref{anexo:pdf1} se puede ver la página pdf generada.


%----------------------------------------------------------------

\subsection{Ejercicio 2}\label{subsec:ej2}
%\subsection{Ejercicio 2}\label{subsec:ej2}

Para que la página esté en horizontal (apaisada) se usa la siguiente instrucción:

\begin{lstlisting}[language=PostScript, caption=Instrucción para página en horizontal, captionpos=b, label=fun:pghorizontal]
<< /PageSize [842 595] >> setpagedevice 
\end{lstlisting}

Para usar distintas tipografías antes del texto se define el tipo de fuente y el tamaño. Por ejemplo:
\begin{lstlisting}[language=PostScript, caption=Tipografía para título, captionpos=b, label=fun:tipografia1]
/Times-Roman findfont % Fuente
25 scalefont % Escala
setfont % Establecer como la fuente actual
\end{lstlisting}
Este es el formato para el texto \textit{PDIH} e \textit{Ines Ruiz}. Para que \textit{Ingeniera Informática} aparezca en cursiva hay que usar la orden:
\begin{lstlisting}[language=PostScript, caption=Tipografía en cursiva, captionpos=b, label=fun:tipografia2]
/Times-Italic findfont 14 scalefont setfont
\end{lstlisting}

Además, a la izquierda se ha dibujado un círculo con el comando \textit{arc}.
\begin{lstlisting}[language=PostScript, caption=Círculo de la tarjeta, captionpos=b, label=fun:circulo]
newpath % Círculo
    0.7 setgray % color gris claro
    3 setlinewidth
    200 300 80 0 360 arc   
    fill
stroke

newpath % Circunferencia
    0.3 setgray % color gris oscuro
    3 setlinewidth
    200 300 80 0 360 arc  
stroke
\end{lstlisting}
Para hacer el efecto de que el borde sea más oscuro que el relleno se han dibujado dos círculos con distintos tonos de grises. En el circulo interno se ha usado la instrucción \textit{fill} para indicar que la figura esté rellena.

A continuación se muestra la imagen obtenida\myref{fig:2ejer2}:
\begin{figure}[H]
\includegraphics[width=0.9\textwidth]{images/ejer2.jpg}
\caption{Propuesta ejercicio 2}
\label{fig:2ejer2}
\end{figure}

El código completo se puede ver en el anexo \ref{anexo:ej2}. En el anexo \ref{anexo:pdf2} se puede ver la página pdf generada.

%----------------------------------------------------------------

\subsection{Ejercicio 3}\label{subsec:ej3}
% \subsection{Ejercicio 3 - Secuencia de cuatro LEDs}\label{subsec:ej3}

Como en el ejercicio anterior, hay que indicar que la página sea apaisada (\myref{fun:pghorizontal}).\\

Para dibujar el corazón se usan dos semicírculos (de 0 hasta 180 grados) y un triángulo (pirámide invertida):
\begin{lstlisting}[language=PostScript, caption=Corazón, captionpos=b, label=fun:corazon]
 % Parte izquierda
newpath
    1 0 0 setrgbcolor  % Color rojo
    3 setlinewidth
    200 300 60 0 180 arc  
    fill
stroke

% Parte derecha
newpath 
    1 0 0 setrgbcolor  % Color rojo
    320 300 60 0 180 arc   
    fill
stroke

% Pico (triángulo)
newpath 
    1 0 0 setrgbcolor  % Color rojo
    140 300 moveto % Esquina superior izquierda
    240 0 rlineto   % Línea hacia la derecha
    -120 -120 rlineto   % Línea hacia abajo-izq
    closepath 
    fill
stroke
\end{lstlisting}

Además, para la poesía se usa la función \textit{show} y con la función \textit{setrgbcolor} se crean distintos tonos de azul jugando con la componente $G$ y $B$.\\

\begin{lstlisting}[language=PostScript, caption=Poesía, captionpos=b, label=fun:poesia]
/Times-Italic findfont 18 scalefont setfont % Fuente
newpath
    0 0.2 1 setrgbcolor  % Azul oscuro
    450 300 moveto 
    (Dale limosna, mujer,) show

    0 0.4 1 setrgbcolor  % Azul medio
    450 280 moveto 
    (que no hay en la vida nada,) show

    0 0.6 1 setrgbcolor  % Azul claro
    450 260 moveto 
    (como la pena de ser) show

    0 0.8 1 setrgbcolor  % Azul muy claro
    450 240 moveto 
    (ciego en Granada.) show

    0 0.1 0.5 setrgbcolor % Contraste para el autor
    470 220 moveto 
    (Francisco de Icaza.) show
stroke
showpage
\end{lstlisting}
A continuación se muestra la imagen obtenida\myref{fig:3ejer}:
\begin{figure}[H]
\includegraphics[width=0.9\textwidth]{images/ejer3.jpg}
\caption{Propuesta ejercicio 3.}
\label{fig:3ejer}
\end{figure}

El código completo se puede ver en el anexo \ref{anexo:ej3}. En el anexo \ref{anexo:pdf3} se puede ver la página pdf generada.

%----------------------------------------------------------------

%% Section 2
\section{Requisitos ampliados}
Como requisitos ampliados (opcionales para subir nota) se propone:
\begin{itemize}
    \item Ejercicio en el que se muestre la imagen de la sonrisa, en apaisado, dentro de una serie de círculos concéntricos de diferentes colores.
    \begin{figure}[H]
        \includegraphics[width=0.7\textwidth]{images/ejemplos_guion/ejemplo4.png}
        \caption{Ejercicio 4. Diana}
        \label{fig:4ejemplo4}
    \end{figure}
    
    \item Ejercicio en el que se creen dos páginas, mostrando en la primera una tabla de datos inventados y en la segunda un texto y una gráfica inventada (los dos ejes y una representación gráfica sencilla en forma de cresta).
    \begin{figure}[H]
        \includegraphics[width=0.7\textwidth]{images/ejemplos_guion/ejemplo5.png}
        \caption{Ejercicio 5. Tabla con gráfica}
        \label{fig:5ejemplo5}
    \end{figure}
\end{itemize}

%----------------------------------------------------------------
\subsection{Ejercicio 4}\label{subsec:ej4}
%\subsection{Ejercicio 4 - Diana}\label{subsec:ej4}
Para hacer los círculos concéntricos de la diana tan solo hay que ir reduciendo el diámetro de los círculos, sin cambiar las coordenadas. Una página A4 apaisada mide $842x595$, por lo que su centro exacto es ($421$, $297$).

\begin{lstlisting}[language=PostScript, caption=Círculos concéntricos, captionpos=b, label=fun:diana]
% Círculo azul (grande)
newpath 
    0.1 0.1 1 setrgbcolor % Azul
    3 setlinewidth
    421 297 180 0 360 arc   
    fill
stroke

% Círculo rojo
newpath 
    1 0 0 setrgbcolor % Rojo
    3 setlinewidth
    421 297 160 0 360 arc   
    fill
stroke

% Círculo verde
newpath 
    0.33 0.67 0 setrgbcolor % Verde
    3 setlinewidth
    421 297 140 0 360 arc   
    fill
stroke

% Círculo lila
newpath 
    1 0.25 0.7 setrgbcolor % Lila
    3 setlinewidth
    421 297 120 0 360 arc   
    fill
stroke

% Círculo amarillo (pequeño)
newpath 
    1 1 0 setrgbcolor % Amarillo
    3 setlinewidth
    421 297 100 0 360 arc   
    fill
stroke
\end{lstlisting}

Para dibujar la cara uso la imagen creada por el profesor en el ejemplo imagen2.ps(\ref{anexo:imagen2}).\\
\begin{lstlisting}[language=PostScript, caption=Imagen dentro de círculos concéntricos, captionpos=b, label=fun:dianaImagen]
370 250 translate
100 100 scale
14 12  8  [14 0 0 -12 0 12]
{<
eeeeeeeeeeeeeeeeeeeeeeee5577
eeeeeeeeeeeeeeeeeeeeeeee5577
eeffffffffffffffffffffee5577
eeffff0000ffff0000ffffee5577
eeffffffffffffffffffffee5577
eeffffffffffffffffffffee5577
eeffff00ffffffff00ffffee5577
eeffffff00000000ffffffee5577
eeeeeeeeeeeeeeeeeeeeeeee5577
eeeeeeeeeeeeeeeeeeeeeeee5577
5577557755775577557755775577
5577557755775577557755775577
>}
image
\end{lstlisting}

A continuación se muestra la imagen obtenida\myref{fig:4ejer}:
\begin{figure}[H]
\includegraphics[width=0.9\textwidth]{images/ejer4.jpg}
\caption{Propuesta ejercicio 4.}
\label{fig:4ejer}
\end{figure}

El programa completo se puede ver en el anexo \ref{anexo:ej4}. En el anexo \ref{anexo:pdf4} se puede ver la página pdf generada.

%----------------------------------------------------------------
\subsection{Ejercicio 5}\label{subsec:ej5}
%\subsection{Ejercicio 5 - Dos páginas}\label{subsec:ej3}

Para que se generen dos páginas se hacen dos llamadas a \textit{showpage}: una al finalizar la tabla y otra al terminar la gráfica.\\

Lo primero es colocar el título, indicando posición y tipografía:
\begin{lstlisting}[language=PostScript, caption=Titulo primera página, captionpos=b, label=fun:titulo]
/Times-Bold findfont 24 scalefont setfont
newpath
    200 700 moveto
    (Mi tabla de datos: Ventas) show
\end{lstlisting}

Para simular la tabla se dibuja una cuadricula:
\begin{lstlisting}[language=PostScript, caption=Cuadricula de la tabla, captionpos=b, label=fun:tabla]
newpath
    % Líneas horizontales
    150 650 moveto 400 650 lineto
    150 620 moveto 400 620 lineto
    150 500 moveto 400 500 lineto
    
    % Líneas verticales
    150 650 moveto 150 500 lineto
    275 650 moveto 275 500 lineto
    400 650 moveto 400 500 lineto
    
    2 setlinewidth 
    0 0 0 setrgbcolor % Negro
stroke
\end{lstlisting}

Con el comando \textit{show} se coloca el texto en la tabla, con el siguiente tipo de letra:
\begin{lstlisting}[language=PostScript, caption=Tipo de fuente, captionpos=b, label=fun:tipo]
/Times-Roman findfont 16 scalefont setfont
\end{lstlisting}

La tabla se rellena con el siguiente contenido:
\begin{lstlisting}[language=PostScript, caption=Relleno de la tabla, captionpos=b, label=fun:relleno]
% Cabeceras
newpath 170 630 moveto (Mes) show
newpath 300 630 moveto (Ventas) show

% Datos inventados
newpath 170 590 moveto (Enero) show
newpath 300 590 moveto (10) show

newpath 170 550 moveto (Febrero) show
newpath 300 550 moveto (50) show

newpath 170 510 moveto (Marzo) show
newpath 300 510 moveto (20) show
\end{lstlisting}

A continuación se muestra la imagen obtenida para la primera página \myref{fig:5ejer1}:
\begin{figure}[H]
    \includegraphics[width=0.9\linewidth]{images/ejer5_pg1.jpg}
    \caption{Propuesta ejercicio 5. Página 1.}
    \label{fig:5ejer1}
\end{figure}

Para dibujar la gráfica primero se pintan los ejes:
\begin{lstlisting}[language=PostScript, caption=Ejes de la gráfica, captionpos=b, label=fun:ejes]
newpath
    % Eje Y (Vertical)
    150 400 moveto
    150 600 lineto 
    
    % Eje X (Horizontal)
    150 400 moveto
    450 400 lineto 
    
    3 setlinewidth 
    0 0 0 setrgbcolor % Negro
stroke
\end{lstlisting}

Dibujo los datos en color rojo:
\begin{lstlisting}[language=PostScript, caption=Gráfica, captionpos=b, label=fun:grafica]
newpath
    150 400 moveto     % Origen
    200 450 lineto     % Sube a 10 (Enero)
    300 580 lineto     % Pico máximo de 50 (Febrero)
    400 480 lineto     % Baja a 20 (Marzo)
    
    2 setlinewidth
    1 0 0 setrgbcolor  % Color rojo
stroke
\end{lstlisting}

Por último, dibujar las etiquetas de los ejes:
\begin{lstlisting}[language=PostScript, caption=Etiquetas de la gráfica, captionpos=b, label=fun:etiq_grafica]
/Times-Roman findfont 12 scalefont setfont
0 0 0 setrgbcolor      % Volver a negro para el texto

% Etiquetas Eje X
newpath 190 380 moveto (Ene) show
newpath 290 380 moveto (Feb) show
newpath 390 380 moveto (Mar) show

% Etiquetas Eje Y
newpath 120 450 moveto (10) show
newpath 120 580 moveto (50) show
\end{lstlisting}

A continuación se muestra la imagen obtenida para la segunda página \myref{fig:5ejer2}:
\begin{figure}[H]
\includegraphics[width=0.9\textwidth]{images/ejer5_pg2.jpg}
\caption{Propuesta ejercicio 5. Página 2.}
\label{fig:5ejer2}
\end{figure}

El programa completo se puede ver en el anexo \ref{anexo:ej5}. En el anexo \ref{anexo:pdf5} se puede ver la página pdf generada.


%----------------------------------------------------------------
%----------------------------------------------------------------

\newpage
\section{Anexo}\label{anexo:an}
% --- Ejer 1
\subsection{Código del Ejercicio 1.} \label{anexo:ej1} 
\lstinputlisting[language=PostScript]{src/ejer1.ps}

\subsection{Documento PDF del Ejercicio 1.} \label{anexo:pdf1}
Documento PDF generado en el ejercicio 1:
\includepdf[pages=-]{pdf/ejercicio1.pdf}

% --- Ejer 2
\subsection{Código del Ejercicio 2} \label{anexo:ej2} 
\lstinputlisting[language=PostScript]{src/ejer2.ps}

\subsection{Documento PDF del Ejercicio 2.} \label{anexo:pdf2}
Al insertar el documento \textit{.pdf} en \textit{LATEX} se gira automáticamente. Documento PDF generado en el ejercicio 2:
\includepdf[pages=-]{pdf/ejercicio2.pdf}

% --- Ejer 3
\subsection{Código del Ejercicio 3} \label{anexo:ej3} 
\lstinputlisting[language=PostScript]{src/ejer3.ps}

\subsection{Documento PDF del Ejercicio 3.} \label{anexo:pdf3}
Documento PDF generado en el ejercicio 3:
\includepdf[pages=-]{pdf/ejercicio3.pdf}

% --- Ejer 4
\subsection{Código del Ejercicio 4} \label{anexo:ej4} 
\lstinputlisting[language=PostScript]{src/ejer4.ps}

\subsection{Documento PDF del Ejercicio 4.} \label{anexo:pdf4}
Documento PDF generado en el ejercicio 4:
\includepdf[pages=-]{pdf/ejercicio4.pdf}

\subsection{Código del profesor para \textit{imagen2}} \label{anexo:imagen2} 
\lstinputlisting[language=PostScript]{src/imagen2.ps}

\subsection{Documento PDF de \textit{imagen2}.} \label{anexo:pdfi2}
Documento PDF generado con el código proporcionado por el profesor:
\includepdf[pages=-]{pdf/imagen2.pdf}

% --- Ejer 5
\subsection{Código del Ejercicio 5} \label{anexo:ej5} 
% Ejer5
\lstinputlisting[language=PostScript]{src/ejer5.ps}

\subsection{Documento PDF del Ejercicio 5.} \label{anexo:pdf5}
Documento PDF de dos páginas generado en el ejercicio 5:
\includepdf[pages=-]{pdf/ejercicio5.pdf}

%\section{Bibliografía y referencias}
%\printbibliography
%\begin{itemize}
%    \item Material proporcionado por el profesor.
%    \item Para ultrasonido:\\ \url{https://www.geekfactory.mx/tutoriales-arduino/tutorial-sensor-ultrasonico-hc-sr04-con-arduino/} ; \url{https://www.luisllamas.es/medir-distancia-con-arduino-y-sensor-de-ultrasonidos-hc-sr04/}
%\end{itemize}

\end{document}